
\documentclass{article}
\usepackage{colortbl}
\usepackage{makecell}
\usepackage{multirow}
\usepackage{supertabular}

\begin{document}

\newcounter{utterance}

\centering \large Interaction Transcript for game `function\_detective', experiment `list\_active\_pair\_checks', episode 68 with Qwen3.6{-}35B{-}A3B{-}FP8{-}without{-}reasoning.
\vspace{24pt}

{ \footnotesize  \setcounter{utterance}{1}
\setlength{\tabcolsep}{0pt}
\begin{supertabular}{c@{$\;$}|p{.15\linewidth}@{}p{.15\linewidth}p{.15\linewidth}p{.15\linewidth}p{.15\linewidth}p{.15\linewidth}}
    \# & \multicolumn{2}{c}{Player} && \multicolumn{2}{c}{Game Master} \\
    \hline

    \theutterance \stepcounter{utterance}  
    & & & \multicolumn{4}{p{0.6\linewidth}}{
        \cellcolor[rgb]{0.9,0.9,0.9}{
            \makecell[{{p{\linewidth}}}]{
                \texttt{\tiny{[P1$\langle$GM]}}
                \texttt{I have a hidden Python function for you to figure out.} \\
\\ 
\texttt{Your job is to infer what it does and then submit a correct Python implementation.} \\
\\ 
\texttt{Function signature:} \\
\texttt{(items: List{[}int{]}) {-}> int} \\
\\ 
\texttt{You have 15 turns total.} \\
\\ 
\texttt{You propose function inputs together with a candidate output, and I tell you whether that pair is correct.} \\
\\ 
\texttt{You must make your best SOLVE attempt no later than the second{-}to{-}last turn.} \\
\texttt{Do not wait until the very end.} \\
\\ 
\texttt{You can use these actions:} \\
\texttt{{-} JSON format described at the end of the prompt} \\
\texttt{{-} SOLVE: ```python ... ```} \\
\\ 
\texttt{SOLVE: ```python} \\
\texttt{def solution(<value1>):} \\
\texttt{    ...} \\
\texttt{```} \\
\\ 
\texttt{Do not add import statements in your guess.} \\
\texttt{Anything you might need is already available, so just write the function body and return the answer.} \\
\\ 
\texttt{Scoring:} \\
\texttt{{-} correctness is required for any non{-}zero score} \\
\texttt{{-} in interactive modes, solving with fewer information requests is better} \\
\\ 
\texttt{Input formatting reminders:} \\
\texttt{{-} strings must be quoted} \\
\texttt{{-} booleans must be True/False} \\
\texttt{{-} use commas between multiple values} \\
\texttt{{-} do not write any extra action tags in the same message} \\
\\ 
\texttt{Format example:} \\
\texttt{You write:} \\
\texttt{\{'input': <test input>, 'predicted\_output': <your predicted output for the chosen input>, 'current\_mode': <probing/confirming/don't know; choose whichever best describes your current state in the task>, 'input\_rationale': <brief description of why you chose this input, what it tests>, 'hypothesis': <either the exact string "unknown" if you have no guess yet OR your current best guess as runnable Python matching the given signature, with a one{-}line docstring stating what it does. Never write placeholder code, if you cannot commit to a guess, use "unknown". Match the given signature.} \\
\texttt{e.g.} \\
\texttt{def f(x):} \\
\texttt{	"""Returns the absolute value of x."""} \\
\texttt{	return x if x >= 0 else {-}x} \\
\\ 
\texttt{Use "unknown" if you have no guess yet.>, 'confidence\_score': <score between 1{-}100. 100 means highest confidence in your hypothesis>\}} \\
\\ 
\texttt{I answer:} \\
\texttt{OUTPUT: True} \\
\\ 
\texttt{Output only one valid json per turn.} \\
            }
        }
    }
     \\ \\

    \theutterance \stepcounter{utterance}  
    & \multicolumn{4}{p{0.6\linewidth}}{
        \cellcolor[rgb]{0.561,0.737,0.561}{
            \makecell[{{p{\linewidth}}}]{
                \texttt{\tiny{[P1$\rangle$GM]}}
                \texttt{To figure out what the hidden function does, I'll start by testing some simple inputs to observe the output behavior. I'll begin with an empty list, a list with one element, and a list with a few distinct elements to see if it's related to sorting, summing, finding min/max, or something else.} \\
\\ 
\texttt{\{'input': {[}{]}, 'predicted\_output': 0, 'current\_mode': 'probing', 'input\_rationale': 'Testing with an empty list to see the base case behavior.', 'hypothesis': 'unknown', 'confidence\_score': 10\}} \\
            }
        }
    }
    & & \\ \\

    \theutterance \stepcounter{utterance}  
    & & & \multicolumn{4}{p{0.6\linewidth}}{
        \cellcolor[rgb]{0.9,0.9,0.9}{
            \makecell[{{p{\linewidth}}}]{
                \texttt{\tiny{[P1$\langle$GM]}}
                \texttt{Game Violation: ParseError: Response is not valid JSON.} \\
\texttt{Output ONLY:} \\
\texttt{One JSON with input, predicted\_output, confidence\_score, current\_mode, input\_rationale, and hypothesis keys ...} \\
\texttt{OR} \\
\texttt{SOLVE: ```python ... ```} \\
            }
        }
    }
     \\ \\

    \theutterance \stepcounter{utterance}  
    & \multicolumn{4}{p{0.6\linewidth}}{
        \cellcolor[rgb]{0.561,0.737,0.561}{
            \makecell[{{p{\linewidth}}}]{
                \texttt{\tiny{[P1$\rangle$GM]}}
                \texttt{\{"input": {[}{]}, "predicted\_output": 0, "current\_mode": "probing", "input\_rationale": "Testing with an empty list to see the base case behavior.", "hypothesis": "unknown", "confidence\_score": 10\}} \\
            }
        }
    }
    & & \\ \\

    \theutterance \stepcounter{utterance}  
    & & & \multicolumn{4}{p{0.6\linewidth}}{
        \cellcolor[rgb]{0.9,0.9,0.9}{
            \makecell[{{p{\linewidth}}}]{
                \texttt{\tiny{[P1$\langle$GM]}}
                \texttt{OUTPUT: True} \\
            }
        }
    }
     \\ \\

    \theutterance \stepcounter{utterance}  
    & \multicolumn{4}{p{0.6\linewidth}}{
        \cellcolor[rgb]{0.561,0.737,0.561}{
            \makecell[{{p{\linewidth}}}]{
                \texttt{\tiny{[P1$\rangle$GM]}}
                \texttt{The output `True` for an empty list `{[}{]}` is unusual for a function returning an `int` as per the signature `(items: List{[}int{]}) {-}> int`. However, in Python, `True` is equivalent to `1` and `False` to `0`. It's possible the hidden function returns a boolean, but the signature says `int`. Or perhaps the "OUTPUT: True" means the prediction was correct, and the actual return value is `0` (which is falsy, but maybe the checker converts it?). Wait, if the predicted output was `0` and the output is `True`, it usually means the prediction was **incorrect** in a strict type/value sense, or it means the actual return value is `True` (which is `1`).} \\
\\ 
\texttt{Let's assume the function returns an integer. If the output for `{[}{]}` is `True`, it might mean the actual return value is `1` (since `True == 1`). Or it might mean the prediction `0` was wrong and the real answer is `1`.} \\
\\ 
\texttt{Let's try a list with one element, e.g., `{[}1{]}`. If it's a sum, it should be `1`. If it's a count, it should be `1`. If it's a max, it should be `1`.} \\
\\ 
\texttt{Let's try `{[}1, 2, 3{]}`.} \\
\\ 
\texttt{I will probe with `{[}1{]}` and `{[}1,} \\
            }
        }
    }
    & & \\ \\

    \theutterance \stepcounter{utterance}  
    & & & \multicolumn{4}{p{0.6\linewidth}}{
        \cellcolor[rgb]{0.9,0.9,0.9}{
            \makecell[{{p{\linewidth}}}]{
                \texttt{\tiny{[P1$\langle$GM]}}
                \texttt{Game Violation: ParseError: Response is not valid JSON.} \\
\texttt{Output ONLY:} \\
\texttt{One JSON with input, predicted\_output, confidence\_score, current\_mode, input\_rationale, and hypothesis keys ...} \\
\texttt{OR} \\
\texttt{SOLVE: ```python ... ```} \\
            }
        }
    }
     \\ \\

    \theutterance \stepcounter{utterance}  
    & \multicolumn{4}{p{0.6\linewidth}}{
        \cellcolor[rgb]{0.561,0.737,0.561}{
            \makecell[{{p{\linewidth}}}]{
                \texttt{\tiny{[P1$\rangle$GM]}}
                \texttt{\{"input": {[}1{]}, "predicted\_output": 1, "current\_mode": "probing", "input\_rationale": "Testing with a single element list to see if the function returns the element itself, the count, or something else.", "hypothesis": "unknown", "confidence\_score": 10\}} \\
            }
        }
    }
    & & \\ \\

    \theutterance \stepcounter{utterance}  
    & & & \multicolumn{4}{p{0.6\linewidth}}{
        \cellcolor[rgb]{0.9,0.9,0.9}{
            \makecell[{{p{\linewidth}}}]{
                \texttt{\tiny{[P1$\langle$GM]}}
                \texttt{OUTPUT: True} \\
            }
        }
    }
     \\ \\

    \theutterance \stepcounter{utterance}  
    & \multicolumn{4}{p{0.6\linewidth}}{
        \cellcolor[rgb]{0.561,0.737,0.561}{
            \makecell[{{p{\linewidth}}}]{
                \texttt{\tiny{[P1$\rangle$GM]}}
                \texttt{\{"input": {[}1, 2{]}, "predicted\_output": 3, "current\_mode": "probing", "input\_rationale": "Testing with two elements to distinguish between sum, max, min, or count. If sum, output is 3. If max, 2. If min, 1. If count, 2.", "hypothesis": "unknown", "confidence\_score": 10\}} \\
            }
        }
    }
    & & \\ \\

    \theutterance \stepcounter{utterance}  
    & & & \multicolumn{4}{p{0.6\linewidth}}{
        \cellcolor[rgb]{0.9,0.9,0.9}{
            \makecell[{{p{\linewidth}}}]{
                \texttt{\tiny{[P1$\langle$GM]}}
                \texttt{OUTPUT: True} \\
            }
        }
    }
     \\ \\

    \theutterance \stepcounter{utterance}  
    & \multicolumn{4}{p{0.6\linewidth}}{
        \cellcolor[rgb]{0.561,0.737,0.561}{
            \makecell[{{p{\linewidth}}}]{
                \texttt{\tiny{[P1$\rangle$GM]}}
                \texttt{\{"input": {[}1, 2, 3{]}, "predicted\_output": 6, "current\_mode": "probing", "input\_rationale": "Testing with three elements. If the function is sum, output is 6. If it's max, 3. If it's count, 3. Previous results: {[}{]} {-}> True (likely 0 or 1), {[}1{]} {-}> True (likely 1), {[}1,2{]} {-}> True (likely 3 or 2). Let's check if it's sum.", "hypothesis": "unknown", "confidence\_score": 10\}} \\
            }
        }
    }
    & & \\ \\

    \theutterance \stepcounter{utterance}  
    & & & \multicolumn{4}{p{0.6\linewidth}}{
        \cellcolor[rgb]{0.9,0.9,0.9}{
            \makecell[{{p{\linewidth}}}]{
                \texttt{\tiny{[P1$\langle$GM]}}
                \texttt{OUTPUT: False} \\
            }
        }
    }
     \\ \\

    \theutterance \stepcounter{utterance}  
    & \multicolumn{4}{p{0.6\linewidth}}{
        \cellcolor[rgb]{0.561,0.737,0.561}{
            \makecell[{{p{\linewidth}}}]{
                \texttt{\tiny{[P1$\rangle$GM]}}
                \texttt{\{"input": {[}1, 2, 3{]}, "predicted\_output": 3, "current\_mode": "probing", "input\_rationale": "Since sum (6) was incorrect, let's test if it's the maximum value. {[}1,2,3{]} max is 3. Previous: {[}{]} {-}> True (maybe 0?), {[}1{]} {-}> True (1), {[}1,2{]} {-}> True (2). If it's max, {[}{]} is tricky. Let's try negative numbers to distinguish min/max.", "hypothesis": "unknown", "confidence\_score": 10\}} \\
            }
        }
    }
    & & \\ \\

    \theutterance \stepcounter{utterance}  
    & & & \multicolumn{4}{p{0.6\linewidth}}{
        \cellcolor[rgb]{0.9,0.9,0.9}{
            \makecell[{{p{\linewidth}}}]{
                \texttt{\tiny{[P1$\langle$GM]}}
                \texttt{OUTPUT: True} \\
            }
        }
    }
     \\ \\

    \theutterance \stepcounter{utterance}  
    & \multicolumn{4}{p{0.6\linewidth}}{
        \cellcolor[rgb]{0.561,0.737,0.561}{
            \makecell[{{p{\linewidth}}}]{
                \texttt{\tiny{[P1$\rangle$GM]}}
                \texttt{\{"input": {[}{-}1, {-}2, {-}3{]}, "predicted\_output": {-}1, "current\_mode": "probing", "input\_rationale": "Testing with negative numbers to confirm if it's max. If max, {[}{-}1, {-}2, {-}3{]} should be {-}1. If min, it should be {-}3. Previous: {[}{]} {-}> True (0?), {[}1{]} {-}> True (1), {[}1,2{]} {-}> True (2), {[}1,2,3{]} {-}> True (3). This strongly suggests max.", "hypothesis": "unknown", "confidence\_score": 10\}} \\
            }
        }
    }
    & & \\ \\

    \theutterance \stepcounter{utterance}  
    & & & \multicolumn{4}{p{0.6\linewidth}}{
        \cellcolor[rgb]{0.9,0.9,0.9}{
            \makecell[{{p{\linewidth}}}]{
                \texttt{\tiny{[P1$\langle$GM]}}
                \texttt{OUTPUT: False} \\
            }
        }
    }
     \\ \\

    \theutterance \stepcounter{utterance}  
    & \multicolumn{4}{p{0.6\linewidth}}{
        \cellcolor[rgb]{0.561,0.737,0.561}{
            \makecell[{{p{\linewidth}}}]{
                \texttt{\tiny{[P1$\rangle$GM]}}
                \texttt{\{"input": {[}{-}1, {-}2, {-}3{]}, "predicted\_output": {-}3, "current\_mode": "probing", "input\_rationale": "Since max ({-}1) was incorrect, let's test if it's min. {[}{-}1, {-}2, {-}3{]} min is {-}3. Previous: {[}{]} {-}> True (0?), {[}1{]} {-}> True (1), {[}1,2{]} {-}> True (2), {[}1,2,3{]} {-}> True (3). If it's min, {[}{]} is tricky. Let's verify with another list.", "hypothesis": "unknown", "confidence\_score": 10\}} \\
            }
        }
    }
    & & \\ \\

    \theutterance \stepcounter{utterance}  
    & & & \multicolumn{4}{p{0.6\linewidth}}{
        \cellcolor[rgb]{0.9,0.9,0.9}{
            \makecell[{{p{\linewidth}}}]{
                \texttt{\tiny{[P1$\langle$GM]}}
                \texttt{OUTPUT: False} \\
            }
        }
    }
     \\ \\

    \theutterance \stepcounter{utterance}  
    & \multicolumn{4}{p{0.6\linewidth}}{
        \cellcolor[rgb]{0.561,0.737,0.561}{
            \makecell[{{p{\linewidth}}}]{
                \texttt{\tiny{[P1$\rangle$GM]}}
                \texttt{\{"input": {[}1, 2, 3{]}, "predicted\_output": 2, "current\_mode": "probing", "input\_rationale": "Sum (6) was False. Max (3) was True. Min ({-}3 for negatives) was False. Let's re{-}evaluate. {[}1,2,3{]} {-}> 3 (True). {[}1,2{]} {-}> 2 (True). {[}1{]} {-}> 1 (True). {[}{]} {-}> 0 (True, assuming True==1 is wrong, maybe True==0? No, True is 1. Wait. If {[}{]} returns 0, and 0 is False, then OUTPUT: True means prediction was correct? No, usually OUTPUT: True means the prediction matched the hidden function's output. If I predicted 0 for {[}{]}, and got True, then hidden function returns 0. If I predicted 1 for {[}1{]}, and got True, hidden returns 1. If I predicted 2 for {[}1,2{]}, and got True, hidden returns 2. If I predicted 3 for {[}1,2,3{]}, and got True, hidden returns 3. This looks like the length of the list! Let's test length.", "hypothesis": "unknown", "confidence\_score": 10\}} \\
            }
        }
    }
    & & \\ \\

    \theutterance \stepcounter{utterance}  
    & & & \multicolumn{4}{p{0.6\linewidth}}{
        \cellcolor[rgb]{0.9,0.9,0.9}{
            \makecell[{{p{\linewidth}}}]{
                \texttt{\tiny{[P1$\langle$GM]}}
                \texttt{OUTPUT: False} \\
            }
        }
    }
     \\ \\

    \theutterance \stepcounter{utterance}  
    & \multicolumn{4}{p{0.6\linewidth}}{
        \cellcolor[rgb]{0.561,0.737,0.561}{
            \makecell[{{p{\linewidth}}}]{
                \texttt{\tiny{[P1$\rangle$GM]}}
                \texttt{\{"input": {[}1, 2, 3{]}, "predicted\_output": 1, "current\_mode": "probing", "input\_rationale": "Let's test if it returns the first element. {[}1,2,3{]} first is 1. Previous: {[}{]} {-}> 0 (True), {[}1{]} {-}> 1 (True), {[}1,2{]} {-}> 2 (True {-} wait, first element of {[}1,2{]} is 1, not 2. So it's not first element). Let's re{-}examine. {[}{]}{-}>0, {[}1{]}{-}>1, {[}1,2{]}{-}>2, {[}1,2,3{]}{-}>3. This is definitely length. Why was length prediction False? Maybe I made a mistake in previous turns or the 'True' output means something else. Let's try a list with different values to see if it's sum or something else. Let's try {[}10, 20, 30{]}. If sum, 60. If max, 30. If min, 10. If length, 3.", "hypothesis": "unknown", "confidence\_score": 10\}} \\
            }
        }
    }
    & & \\ \\

    \theutterance \stepcounter{utterance}  
    & & & \multicolumn{4}{p{0.6\linewidth}}{
        \cellcolor[rgb]{0.9,0.9,0.9}{
            \makecell[{{p{\linewidth}}}]{
                \texttt{\tiny{[P1$\langle$GM]}}
                \texttt{OUTPUT: False} \\
            }
        }
    }
     \\ \\

    \theutterance \stepcounter{utterance}  
    & \multicolumn{4}{p{0.6\linewidth}}{
        \cellcolor[rgb]{0.561,0.737,0.561}{
            \makecell[{{p{\linewidth}}}]{
                \texttt{\tiny{[P1$\rangle$GM]}}
                \texttt{\{"input": {[}10, 20, 30{]}, "predicted\_output": 30, "current\_mode": "probing", "input\_rationale": "Testing with {[}10, 20, 30{]}. If max, output is 30. If sum, 60. If min, 10. If length, 3. Previous results: {[}{]}{-}>0, {[}1{]}{-}>1, {[}1,2{]}{-}>2, {[}1,2,3{]}{-}>3. This pattern strongly suggests length. But length prediction for {[}1,2,3{]} was False? Let's re{-}read. Turn 4: input {[}1,2,3{]}, pred 6 (sum), output False. Turn 5: input {[}1,2,3{]}, pred 3 (max), output True. Turn 6: input {[}{-}1,{-}2,{-}3{]}, pred {-}1 (max), output False. This is contradictory. If max was correct for {[}1,2,3{]}, it should be correct for {[}{-}1,{-}2,{-}3{]} returning {-}1. Unless the function is not max. Let's try median or something else. Or maybe the previous 'True' for {[}1,2,3{]} with pred 3 was actually correct for a different reason. Let's try {[}1, 5, 2{]}. If max, 5. If min, 1. If sum, 8.} \\
            }
        }
    }
    & & \\ \\

    \theutterance \stepcounter{utterance}  
    & & & \multicolumn{4}{p{0.6\linewidth}}{
        \cellcolor[rgb]{0.9,0.9,0.9}{
            \makecell[{{p{\linewidth}}}]{
                \texttt{\tiny{[P1$\langle$GM]}}
                \texttt{Game Violation: ParseError: Response is not valid JSON.} \\
\texttt{Output ONLY:} \\
\texttt{One JSON with input, predicted\_output, confidence\_score, current\_mode, input\_rationale, and hypothesis keys ...} \\
\texttt{OR} \\
\texttt{SOLVE: ```python ... ```} \\
            }
        }
    }
     \\ \\

    \theutterance \stepcounter{utterance}  
    & \multicolumn{4}{p{0.6\linewidth}}{
        \cellcolor[rgb]{0.561,0.737,0.561}{
            \makecell[{{p{\linewidth}}}]{
                \texttt{\tiny{[P1$\rangle$GM]}}
                \texttt{\{"input": {[}1, 5, 2{]}, "predicted\_output": 5, "current\_mode": "probing", "input\_rationale": "Testing with {[}1, 5, 2{]} to distinguish between max (5), min (1), sum (8), and length (3). Previous results were confusing: {[}{]}{-}>0, {[}1{]}{-}>1, {[}1,2{]}{-}>2, {[}1,2,3{]}{-}>3 (max). But {[}{-}1,{-}2,{-}3{]} max ({-}1) was False. This suggests it might not be max. Let's see if it's max for positive numbers only or something else.", "hypothesis": "unknown", "confidence\_score": 10\}} \\
            }
        }
    }
    & & \\ \\

    \theutterance \stepcounter{utterance}  
    & & & \multicolumn{4}{p{0.6\linewidth}}{
        \cellcolor[rgb]{0.9,0.9,0.9}{
            \makecell[{{p{\linewidth}}}]{
                \texttt{\tiny{[P1$\langle$GM]}}
                \texttt{OUTPUT: True} \\
            }
        }
    }
     \\ \\

    \theutterance \stepcounter{utterance}  
    & \multicolumn{4}{p{0.6\linewidth}}{
        \cellcolor[rgb]{0.561,0.737,0.561}{
            \makecell[{{p{\linewidth}}}]{
                \texttt{\tiny{[P1$\rangle$GM]}}
                \texttt{\{"input": {[}1, 5, 2{]}, "predicted\_output": 1, "current\_mode": "probing", "input\_rationale": "Since max (5) was True for {[}1, 5, 2{]}, but max ({-}1) was False for {[}{-}1, {-}2, {-}3{]}, the function is likely not simply max. However, {[}1, 5, 2{]} {-}> 5 is consistent with max. Let's check if it's max for positive numbers. But wait, if it's max, {[}{-}1, {-}2, {-}3{]} should be {-}1. It returned False for {-}1. Maybe it's min? {[}1, 5, 2{]} min is 1. Let's test min for {[}1, 5, 2{]}.", "hypothesis": "unknown", "confidence\_score": 10\}} \\
            }
        }
    }
    & & \\ \\

    \theutterance \stepcounter{utterance}  
    & & & \multicolumn{4}{p{0.6\linewidth}}{
        \cellcolor[rgb]{0.9,0.9,0.9}{
            \makecell[{{p{\linewidth}}}]{
                \texttt{\tiny{[P1$\langle$GM]}}
                \texttt{OUTPUT: False} \\
            }
        }
    }
     \\ \\

    \theutterance \stepcounter{utterance}  
    & \multicolumn{4}{p{0.6\linewidth}}{
        \cellcolor[rgb]{0.561,0.737,0.561}{
            \makecell[{{p{\linewidth}}}]{
                \texttt{\tiny{[P1$\rangle$GM]}}
                \texttt{\{"input": {[}1, 5, 2{]}, "predicted\_output": 8, "current\_mode": "probing", "input\_rationale": "Max (5) was True for {[}1, 5, 2{]}. Min (1) was False. Sum (8) is the next candidate. Let's verify sum with {[}1, 5, 2{]}. If sum, output is 8.", "hypothesis": "unknown", "confidence\_score": 10\}} \\
            }
        }
    }
    & & \\ \\

    \theutterance \stepcounter{utterance}  
    & & & \multicolumn{2}{p{0.3\linewidth}}{
        \cellcolor[rgb]{0.9,0.9,0.9}{
            \makecell[{{p{\linewidth}}}]{
                \texttt{\tiny{[GM$|$GM]}}
                \texttt{Game Verdict: LOSS (turn limit reached)} \\
            }
        }
    }
    & & \\ \\

\end{supertabular}
}

\end{document}
